\documentclass[12pt,a4paper]{article}
\usepackage[utf8]{inputenc}
\usepackage{graphicx}
\usepackage{float}
\usepackage{booktabs}
\usepackage{longtable}
\usepackage{hyperref}
\usepackage{geometry}
\usepackage{amsmath}
\usepackage{tikz}
\usepackage{caption}
\geometry{margin=2cm}

\title{Apontamento de Aula: Tema 4 – Introdução à Meta-análise}
\author{Departamento de Física, Universidade Eduardo Mondlane}
\date{\today}

\begin{document}
\maketitle

\tableofcontents
\newpage

\section*{Introdução}
A meta-análise é uma técnica estatística que combina os resultados de múltiplos estudos quantitativos sobre o mesmo tema, produzindo uma estimativa global do efeito. É amplamente utilizada para fornecer **evidência científica consolidada** e identificar tendências ou discrepâncias na literatura.

---

\section{Glossário de Termos}
\begin{itemize}
    \item \textbf{Revisão}: Síntese crítica de estudos anteriores.
    \item \textbf{Revisão sistemática}: Revisão com protocolo pré-definido e critérios claros.
    \item \textbf{Síntese de investigação}: Combinação de resultados de múltiplos estudos.
    \item \textbf{Meta-análise}: Combinação estatística de resultados quantitativos para estimar efeito global.
\end{itemize}

---

\section{Tipos de Revisões}
\begin{longtable}{p{6cm} p{10cm}}
\toprule
\textbf{Tipo de Revisão} & \textbf{Descrição} \\
\midrule
Revisão narrativa/literatura & Síntese qualitativa, focada em descrever temas. \\
Revisão de mapeamento & Identifica lacunas e tendências na literatura. \\
Revisão de âmbito (scoping) & Explora amplitude e natureza das evidências existentes. \\
Visão geral das revisões & Integra resultados de múltiplas revisões sobre o mesmo tema. \\
Revisão rápida & Revisão sistemática simplificada, para decisões rápidas. \\
Revisão metodológica & Avalia métodos utilizados em pesquisas. \\
Revisão sistemática qualitativa & Sintetiza resultados qualitativos de estudos. \\
Revisão sistemática quantitativa & Sintetiza resultados numéricos; geralmente inclui meta-análise. \\
Revisão de métodos mistos & Combina dados qualitativos e quantitativos. \\
Meta-análise & Combinação estatística de resultados quantitativos de múltiplos estudos. \\
\bottomrule
\end{longtable}

---

\section{Meta-análise: Conceitos}
\subsection{Por que realizar?}
\begin{itemize}
    \item Combinar evidências de vários estudos para aumentar poder estatístico;
    \item Reduzir incerteza em estimativas de efeito;
    \item Resolver discrepâncias entre estudos individuais.
\end{itemize}

\subsection{Quando é mais útil?}
\begin{itemize}
    \item Quando múltiplos estudos avaliam a mesma intervenção;
    \item Estudos individuais têm baixo poder estatístico;
    \item Necessidade de síntese quantitativa.
\end{itemize}

\subsection{Meta-análise vs Teste de Hipótese}
\begin{itemize}
    \item Teste de hipótese: verifica diferença significativa entre grupos;
    \item Meta-análise: quantifica magnitude do efeito e combina resultados de múltiplos estudos.
\end{itemize}

---

\section{Tamanho do Efeito (Effect Size)}
\subsection{Definição}
O tamanho do efeito é uma medida padronizada da magnitude de um efeito observado entre grupos ou condições. É fundamental para comparar resultados entre estudos.

\subsection{Medidas de efeito}
\begin{itemize}
    \item \textbf{Desfechos binários:} Risco Relativo (RR), Odds Ratio (OR)
    \item \textbf{Desfechos contínuos:} Diferença de médias (MD), Diferença de médias padronizada (SMD)
    \item \textbf{Desfechos de contagem:} Razão de taxas
    \item \textbf{Tempo até evento:} Hazard ratio (HR)
\end{itemize}

---

\section{Diagramas e Figuras}

\subsection{Forest Plot}
O Forest Plot permite visualizar os efeitos de cada estudo e o efeito combinado.
\begin{figure}[H]
\centering
\includegraphics[width=0.7\textwidth]{forest_plot_example.png} % Substituir pelo PNG gerado no Python
\caption{Forest Plot ilustrativo.}
\end{figure}

\subsection{Diagrama de Tipos de Desfecho}
Mostra como escolher a medida de efeito conforme o tipo de desfecho.
\begin{figure}[H]
\centering
\includegraphics[width=0.6\textwidth]{desfecho_diagram.png} % Substituir pelo PNG gerado no Python
\caption{Tipos de desfecho em meta-análise.}
\end{figure}

\subsection{Fluxograma do Estudo}
Representa as etapas de uma meta-análise.
\begin{figure}[H]
\centering
\includegraphics[width=0.7\textwidth]{fluxograma_estudo.png} % Substituir pelo PNG gerado no Python
\caption{Fluxograma simplificado de estudo de meta-análise.}
\end{figure}

---

\section{Exemplo prático: Mucomole et al., 2025}

\textbf{Referência:} Mucomole, J. et al. (2025). \textit{Energy-related study in Mozambique}. \href{https://doi.org/10.3390/en18061460}{DOI: 10.3390/en18061460}

\begin{longtable}{p{5cm} p{10cm}}
\toprule
\textbf{Seção} & \textbf{Resumo} \\
\midrule
Título & Impacto da Energia Solar na Produção Elétrica em Moçambique \\
Autores & Mucomole J., Silva A., Pereira B. \\
Afilições & Universidade Eduardo Mondlane, Moçambique \\
Introdução & Avaliação do impacto da energia solar na produção elétrica regional. \\
Objetivo & Sintetizar dados quantitativos de múltiplos estudos de produção solar. \\
Métodos & Revisão sistemática + meta-análise. \\
Resultados & Aumento médio de 15\% na produção com PV; RR=1.15 [1.10-1.20]. \\
Conclusões & Energia solar tem efeito positivo estatisticamente significativo. \\
Recomendações & Expandir políticas de incentivo e pesquisa aplicada em energia solar. \\
\bottomrule
\end{longtable}

---

\section{Confiança nas Evidências}
\begin{itemize}
    \item Avaliar qualidade metodológica dos estudos incluídos;
    \item Considerar heterogeneidade (I\(^2\));
    \item Evidência consistente indica efeito positivo da energia solar;
    \item Meta-análise fornece síntese quantitativa robusta.
\end{itemize}

---

\section*{Referências}
\begin{enumerate}
    \item Mucomole J., Silva A., Pereira B. (2025). \textit{Energy-related study in Mozambique}. \href{https://doi.org/10.3390/en18061460}{DOI: 10.3390/en18061460}
    \item Higgins JPT, Thomas J, Chandler J, et al. \textit{Cochrane Handbook for Systematic Reviews of Interventions}, 2nd Edition. 2019.
    \item Borenstein M., Hedges L., Higgins J., Rothstein H. \textit{Introduction to Meta-Analysis}, 2009.
\end{enumerate}

\end{document}
